\section{Discussion}
\label{sec:discussion}

\subsection{Why these two issues matter beyond transit}

The two issues isolated in \Cref{sec:problem:failures}---weak action-conditional signal under sparse actions, and reward-scale drift under multi-fidelity training---are not specific to bus holding.
Both arise from structural properties that appear in other multi-fidelity RL domains.

\paragraph{Sparse-action domains}
Any setting where the action space is low-dimensional or nearly discrete, and where inter-decision dynamics are dominated by exogenous factors, exhibits a small $I(S'; A \mid S)$.
Examples beyond transit include ramp metering, signal phase selection, crew dispatching, and demand-responsive transit---all problems where an agent makes one decision per event and the intervening dynamics are governed by traffic and passengers.
In such domains, a DARC-style factored discriminator's access to an action-specific signal is limited.
The contrastive alternative, conditioning only on $(s, s')$, is directly applicable subject to \Cref{ass:action_invariance} holding on domain data.

\paragraph{Multi-fidelity simulators with reward-magnitude gaps}
Reward-scale differences between fidelity levels are common: simplified contact models in robotics produce smaller effort penalties, simplified demand models in power systems produce smaller headroom-violation penalties, and so on.
Whenever the simulator and evaluation environments report rewards on different magnitudes, an LCB-only sim-side target drifts toward the simulator's scale over repeated updates.
A one-sided bootstrap Q-floor offers a cheap fix that does not require pre-computed MC returns or reward statistics, but it is not the only option---explicit reward rescaling or learned reward calibration are equally reasonable first-line treatments (and we evaluate a rescaling baseline in \Cref{sec:experiments}).

\subsection[Calibration of the high-fidelity simulator against the real corridor]{Calibration of $\Mreal$ against the real corridor}
\label{sec:disc:calibration}

We use ``real corridor'' to refer to the topology, schedule, and bus fleet of an existing 12-line, 389-bus high-frequency operating area; the SUMO microsimulation $\Mreal$ replays this topology under a fixed passenger origin--destination (OD) profile and a single deterministic traffic seed.
This is sufficient to establish that the framework produces meaningful behaviour on a realistic-scale network, but it is \emph{not} a calibration claim: route-level travel-time error against AVL/APC observations, dwell-time fitting against IC-card data, OD validation against passenger-survey or AFC records, and train/calibration/test day separation are all out of scope for this version.
Future validation should report (i) route-level travel-time MAPE and headway-distribution overlap against an aligned real operating day, (ii) dwell-time fitting error, (iii) bunching-rate and headway-CV agreement between $\Mreal$ and the observed corridor, and (iv) explicit calibration vs.\ test-day separation.
Until those validations are available, results in this paper should be read as evidence on a real-topology synthetic-demand benchmark rather than on a calibrated digital twin.

\subsection{Remaining evaluation-scope gaps}
\label{sec:disc:eval_scope}

The multi-scenario evaluation in \Cref{sec:exp:multiseed} (Tab.~\ref{tab:multiseed}) covers $3$ SUMO seeds $\times$ $3$ OD scales = $9$ scenarios per method under common random numbers, addressing the basic generalisation question.
Three gaps remain.
First, the OD profile is a single workday template scaled multiplicatively; we have not evaluated weekend OD, special-event peaks, or empirically observed day-to-day variation drawn from real ridership records.
Second, no incident scenarios (link blockage, lane closure, signal failure) are included; the policy's robustness to non-stationary disturbances is therefore unknown.
Third, there is no train/calibration/test day separation: all $9$ scenarios share the same underlying network model, with only seed and demand-scale perturbation.
These extensions are consistent with the calibration roadmap in \Cref{sec:disc:calibration} and should be reported jointly when AVL/APC-validated $\Mreal$ becomes available.

\subsection{Reward shape vs operational utility}
\label{sec:disc:operational}

The composite reward used during training---a weighted sum of headway-deviation magnitude, holding penalty, and forward-looking headway error---is a learning-side proxy, not a direct measure of operator value.
\Cref{tab:multiseed} surfaces a concrete instance: $200$-epoch pure-online SAC has the worst cumulative reward among the four RL rows reported but the lowest mean passenger waiting time and lowest headway CV; the $1{,}000$-epoch pure-online SAC retrain trades some of that waiting-time advantage for higher reward, sitting closer to H2O+ on the reward axis but with passenger waiting closer to the rest.
The mechanism is straightforward: the per-second holding penalty dominates the future bunching reduction the hold buys, so policies that hold less are reward-favoured even when they leave passengers waiting marginally longer.
For deployment this means cumulative reward should be read as a proxy, not a target; the passenger-side columns of \Cref{tab:multiseed} and \Cref{tab:passenger_full} are diagnostic comparison points that still require calibration against observed passenger records before they can support an operational-benefit claim.
We do not retune the reward weights here---that would change the optimisation landscape across the entire ablation matrix---but flag this as a calibration step required before any operator pilot.

\subsection{Limitations}

\paragraph{Action-conditional domain invariance may not hold elsewhere}
\Cref{ass:action_invariance} requires that the simulator and real environment share action semantics up to the gap in state-transition distributions.
On our benchmark the data-level test reports $\hat{r} = 0.156$ (\Cref{sec:exp:analysis}), well below the working $0.30$ threshold we adopt; the contrastive discriminator's IS-weight estimates are therefore plausibly grounded on this benchmark, though the threshold itself is a heuristic and the test's coverage is limited to state regions densely represented in $\Doff$.
On a different benchmark---e.g.\ a sim that ignores hold commands while the real system honours them---the assumption can fail, in which case the contrastive variant should be replaced by an action-conditioned discriminator (TransDisc or DynamicsDisc).
The data-level test in \Cref{sec:exp:analysis} (and the script \texttt{verify\_assumption1.py}) is the recommended pre-deployment check.

\paragraph{Data composition sensitivity}
The offline-Q floor is only as informative as the offline dataset is diverse.
Our $\Doff$ combines four behaviour policies (SAC, heuristic, $\epsilon$-random, RE-SAC SUMO expert ep39); homogeneous offline data from a single policy would yield a narrow $q_{\text{floor}}$ that may not span the Q-scale actually encountered during online training.
For practitioners, this suggests deliberately diversifying offline data collection where possible.

\paragraph{The offline-Q floor is an engineering fix, not a theorem}
The bootstrap floor in \Cref{eq:qfloor} is a one-sided, scalar-level anchor; it does not eliminate all effects of reward-scale mismatch.
A policy that is optimal in $\Msim$ and slightly suboptimal in $\Mreal$ can still be preferred under the hybrid loss when the sim gradient dominates.
Near-optimal policy transfer likely still requires either explicit reward calibration \citep{hanna2017grounded} or curriculum rollout schedules, neither of which this paper studies.

\paragraph{Single-corridor evaluation}
Our benchmark is a single 12-line corridor.
Although the offline dataset is large and the evaluation episode is long, generalisation across corridors---networks of different size, demand patterns, or operating agencies---is not tested here.
We discuss this as the main direction for replication below.

\subsection{Deployment considerations}

The framework is designed with practical deployment in mind: offline data is cheaply collected from existing operational records, online interactions occur in a simulator under the agency's control, and no policy update requires touching live fleets.
Nonetheless, a responsible deployment pipeline should include at least three additional gates.
First, policy outputs should be reviewed against operational constraints (e.g., maximum allowed hold time per stop) before activation.
Second, the trained policy should be first evaluated on held-out historical operating days in the SUMO-like environment, not on live passengers.
Third, rollout should begin on a single route or off-peak window and expand gradually based on measured service indicators.
None of these precautions are unique to our method, but all are necessary when deploying any RL policy in a real transit network.

\subsection{Toward true sim-to-real and open questions}

This paper deliberately studies a \emph{multi-fidelity sim-to-sim} problem; establishing sim-to-real transfer requires two further steps that we identify but do not undertake.
First, historical-data replay: given operator AVL/GPS records, evaluate the trained policy by replaying observed arrival events and comparing simulated vs observed holding outcomes.
Second, gated pilot deployment: instrument a single route or off-peak window with supervisor approval, starting with constrained action ranges and expanding based on measured service indicators.
Both steps face institutional and data-access constraints orthogonal to the methodological questions addressed here.

Beyond transit, three extensions warrant investigation.
Integrating the offline-Q floor with \emph{distributional} RL (e.g., IQN, DSAC) would give explicit risk-sensitive quantile anchors.
Extending to three or more simulators of increasing fidelity---a true \emph{multi-fidelity hierarchy}---would allow cheap rollouts for exploration and expensive rollouts for critical evaluation, with the discriminator and Q-floor applied pairwise.
Joint timetable-and-holding optimisation under a multi-fidelity dynamics gap is a natural next problem already of operational interest; it would require extending the current event-level holding policy to a higher-level timetable or dispatch decision layer.
